% TY61000C
% Wundversorgung
% 14.0
% 2013-11-30
\label{m:wundversorgung}
\begin{itemize}
  \item Wunden dürfen nicht direkt berührt werden, außer es
  ist zur raschen Blutstillung unbedingt notwendig! 
Es
  muss immer mit Handschuhen und mit \EE{sterilem Material}
  gearbeitet werden!\footnote{Bei der Wundversorgung sind
  immer sterile Materialien zu verwenden. Papierhandtücher
  o.\,ä. sind \E{nicht zulässig}! Eine Ausnahme kann
  lediglich bei der Versorgung des Stichs bei der
  Blutzuckermessung mittels eines sauberen Tupfers gemacht
  werden, da die Wunde minimal klein ist.}
  
  % \item Wunden sollen grundsätzliche \E{nicht} gereinigt
  % werden! Sie dürfen nur lose Glassplitter oder ähnliches
  % vorsichtig aus der Wunde entfernen.
  \item \EE{Reinigung}: Oberflächliches \EE{ab}spülen mit
  steriler \I[Kochsalzlösung!pyhsiologische]{physiologischer
  Kochsalzlösung} (\ce{NaCl}
  0,9\PC*) von innen
  nach außen. Bei oberflächlichen
  Wunden ist auch ein Abspülen unter fließendem \IX{Leitungswasser} 
  zulässig \cite{Fernandez2012}. 
  Wunden, die im Spital
  weiter versorgt werden, sollen nicht übermäßig gereinigt werden, 
  auch soll keine
  wesentliche Zeitverzögerung entstehen.  
  \Commentary[Wundreinigung]{Alternativ sind
  statt der physiologischen Kochsalzlösung auch andere
  kristalloide Infusionslösungen verwendbar. Die
  Desinfektion mittels Wunddesinfektionsmittels ist nicht
  mehr empfohlen. \cite{Fernandez2012}}
  \item \EE{Fremdkörper} (Messer, Schere oder sonstige
  Pfählungsgegenstände) sind \EE{in der Wunde zu belassen}!
  Diese Gegenstände sind mit geeignetem Fixationsmaterial
  (z.\,B. Mullbinden) zu stabilisieren. Lose, kleine
  Fremdkörper wie z.\,B. Glassplitter dürfen mittels \E{steriler Instrumente} (Pinzette, {\ldots}) entfernt
  werden.
  \item Ausgetretene Strukturen (z.\,B. Gehirn, Eingeweide)
  dürfen \E{nicht zurückgestopft} werden! Solche Organteile
  sollten mit steriler physiologischer Kochsalzlösung
feucht gehalten und anschließend steril abgedeckt werden! 
  \item Ein steriler
  Wundverband mit der jeweiligen Situation angemessenem
  Material (z.\,B. Pflaster, Momentverband, Kompresse mit
  Mullbinde, Wundfolie, \ldots) muss angelegt oder die
  Wunde anderweitig steril, aber \E{nicht} luftdicht
  abgedeckt werden (Ausnahme: Bagatellverletzungen).
\end{itemize}